\section{The linear criterion}\label{sec:criterion}

Identify the hold with the zero vector of $\R^m$, so that $D(u,\theta)=2\langle u,\theta\rangle-5$ for every action $u\in U_m\subset\{-1,0,1\}^m$. For a nonempty $S\subseteq\Theta_m$ write $\Delta(S)$ for the probabilities on $S$ and
\begin{equation}\label{eq:V}
V(S)\;=\;\max_{y\in[-1,1]^m}\ \min_{\theta\in S}\ \langle\theta,y\rangle .
\end{equation}

\begin{theorem}[linear criterion]\label{thm:criterion}
For a nonempty set $S$ of cells the following are equivalent.
\begin{enumerate}[label=(\roman*),leftmargin=2.2em,itemsep=1pt]
\item $S$ is blind-survivable from some finite slack.
\item $V(S)\ge\tfr52$: there is a mixed action $y\in[-1,1]^m$, which may be taken rational, with $\langle\theta,y\rangle\ge\tfr52$ for every $\theta\in S$.
\item $\min_{q\in\Delta(S)}\big\|\sum_{\theta\in S}q_\theta\theta\big\|_1\ \ge\ \tfr52$.
\end{enumerate}
Moreover $V(S)=\min_{q\in\Delta(S)}\|\sum_\theta q_\theta\theta\|_1$, a minimiser $q$ can be chosen with $|\supp q|\le m+1$, and when $V(S)\ge\tfr52$ a maximiser $y$ of \eqref{eq:V} can be chosen rational with denominators at most those of a vertex of the polytope $\{y\in[-1,1]^m:\langle\theta,y\rangle\ge V(S)\ \forall\theta\in S\}$.
\end{theorem}

\begin{proof}
(i)$\Rightarrow$(ii). Let the word $(u_t)$ keep $S$ from slack $L$, and put $\bar y_T=\tfr1T\sum_{t<T}u_t\in[-1,1]^m$. Survival at time $T$ reads $L+\sum_{t<T}(2\langle u_t,\theta\rangle-5)\ge0$, i.e. $2\langle\bar y_T,\theta\rangle-5\ge-L/T$ for every $\theta\in S$. A limit point $y$ of $(\bar y_T)$ in the compact cube satisfies $\langle y,\theta\rangle\ge\tfr52$ for all $\theta\in S$. The feasible set $\{y\in[-1,1]^m:\langle\theta,y\rangle\ge\tfr52\ \forall\theta\in S\}$ is a nonempty rational polytope, so it has a rational vertex.

(ii)$\Rightarrow$(i). Let $y$ be rational with $\langle\theta,y\rangle\ge\tfr52$ on $S$. The product measure $p(u)=\prod_{i}\tfr12(1+u_iy_i)$ on $\{\pm1\}^m$ has mean $y$. Let $N$ be a common denominator of the $p(u)$ and let $w$ be the word of length $N$ that plays each $u$ exactly $Np(u)$ times, repeated periodically. Its total drift over one period is $\sum_uNp(u)D(u,\theta)=N(2\langle y,\theta\rangle-5)\ge0$ for every $\theta\in S$, so the slack of every cell at the end of each period is at least its value at the start, and the slack along the word never falls below $-\delta$, where $\delta=\max_{\theta\in S}\max_{t\le N}\big(-\sum_{\tau<t}D(w_\tau,\theta)\big)^+$ is the dip of the first period. Hence $S$ is blind-survivable from slack $\delta$.

(ii)$\Leftrightarrow$(iii) and the value identity. The payoff $(y,q)\mapsto\langle y,\sum_\theta q_\theta\theta\rangle$ is bilinear on the product of the compact convex sets $[-1,1]^m$ and $\Delta(S)$, so by von Neumann's minimax theorem \cite{vonNeumann1928}
\[V(S)=\max_y\min_{\theta\in S}\langle\theta,y\rangle=\max_y\min_q\Big\langle y,\sum_\theta q_\theta\theta\Big\rangle=\min_q\max_y\Big\langle y,\sum_\theta q_\theta\theta\Big\rangle=\min_q\Big\|\sum_\theta q_\theta\theta\Big\|_1 .\]
The support bound: \eqref{eq:V} is the linear programme $\max\{t:\ t-\langle\theta,y\rangle\le0\ (\theta\in S),\ -1\le y_i\le1\}$ in the $m+1$ variables $(y,t)$; a basic optimal solution of its dual has at most $m+1$ nonzero multipliers, and the multipliers of the rows indexed by $S$, normalised, are a minimiser $q$ (complementary slackness gives $\sum_\theta q_\theta=1$ because $t$ is free).
\end{proof}

The two sides of the theorem are \emph{certificates}. A rational $y$ is checked by $|S|$ inner products; a rational $q$ by one $\ell_1$ norm. The exact module computes both by a small simplex in rational arithmetic and confirms, for every set it classifies, that $\min_{\theta\in S}\langle\theta,y\rangle=\|\sum_\theta q_\theta\theta\|_1$ exactly. Because the data are integers, every $V(S)$ is a rational with small denominator; the values met in this paper are $0,1,2,\tfr52,3,\dots$, and the boundary value $\tfr52$ occurs (Section~\ref{sec:classification}).

\begin{corollary}[finite horizon of a non-survivable set]\label{cor:horizon}
Let $V(S)<\tfr52$ and let $q$ be a minimiser in (iii). From common slack $L$, every blind word drives some cell of $\supp q$ below the floor within $\lfloor L/(5-2V(S))\rfloor+1$ steps; equivalently no word keeps $S$ for more than $\lfloor L/(5-2V(S))\rfloor$ steps.
\end{corollary}
\begin{proof}
If a word keeps $S$ for $T$ steps, the averaging of the previous proof gives $2\langle\bar y_T,\theta\rangle-5\ge-L/T$ for $\theta\in\supp q$; pooling with $q$ and using $\langle\bar y_T,\sum_\theta q_\theta\theta\rangle\le\|\sum_\theta q_\theta\theta\|_1=V(S)$ yields $2V(S)-5\ge-L/T$, i.e. $T\le L/(5-2V(S))$.
\end{proof}

\begin{corollary}[pairs]\label{cor:pairs}
For two cells at Hamming distance $h$, $V(\{a,b\})=m-h$; the pair is survivable iff $h\le m-\tfr52$, i.e. $h\le\lfloor m-\tfr52\rfloor=:\delta_m$ ($\delta_4=1,\delta_5=2,\delta_6=3$). A pair at distance $h>\delta_m$ survives at most $\lfloor L/(5-2(m-h))\rfloor$ steps; for $h=\delta_m+1$ this is $L$ steps, attained for even $m$. The whole cube, and any set containing a cell and its antipode, is never blind-survivable.
\end{corollary}
\begin{proof}
With $y_i=a_i$ on the $m-h$ agreeing coordinates and $y_i=0$ elsewhere, $\langle y,a\rangle=\langle y,b\rangle=m-h$; with $q=(\tfr12,\tfr12)$, $\|\tfr12(a+b)\|_1=m-h$. For $h=\delta_m+1$ one has $m-h=2$ for every $m$, so the bound is $L$ steps. For even $m$ the constant action at distance $\tfr{m-2}2$ from both cells costs each exactly one tenth per step and attains it; for $m=5$ no action is equidistant from the two cells and the exact horizon, by forward search, is $2\lfloor(L-1)/2\rfloor$ for $L\le20$ (Table~S2.2). The antipode of $a$ has $h=m$ and $V=0$.
\end{proof}

The second statement is the \emph{pair-sum lemma} of the earlier editions in its sharp form: for a pair, $D(u,a)+D(u,b)=4\langle u,\tfr12(a+b)\rangle-10\le4(m-h)-10$, with equality for a matched action, so the pair's total slack is non-increasing exactly when $h>m-\tfr52$. The far-pair horizon $10(z_0-1)$ of the Lean module \texttt{EBC\_Pairs\_v3} (\texttt{far\_pair\_horizon\_bound}) is the case $h=\delta_4+1=2$ of Corollary~\ref{cor:horizon}.

\begin{corollary}[the floor and Hamming balls]\label{cor:balls}
$S$ is blind-survivable at the floor ($L=0$) iff $S\subseteq\Ball(u,r^\ast(m))$ for some sign action $u$, where $r^\ast(m)=\lfloor(2m-5)/4\rfloor$. Hence $\calM_m(0)$ consists of the $2^m$ Hamming balls of radius $r^\ast(m)$, of sizes $1,6,7,29,37$ for $m=4,\dots,8$.
\end{corollary}
\begin{proof}
From zero slack every step needs a nonnegative drift for every cell, so every action used satisfies $S\subseteq\{\theta:D(u,\theta)\ge0\}=\Ball(u,r^\ast)$; conversely the constant action $u$ keeps $\Ball(u,r^\ast)$. Balls of radius $r^\ast$ are maximal among balls and, having a common centre action, pairwise incomparable.
\end{proof}

\begin{corollary}[diameter and size]\label{cor:kleitman}
A survivable set has Hamming diameter at most $\delta_m$. By Kleitman's theorem \cite{Kleitman1966} its size is at most $K(m)=\sum_{i\le r}\binom mi$ if $\delta_m=2r$ and $2\sum_{i\le r}\binom{m-1}i$ if $\delta_m=2r+1$: $K(4)=2$, $K(5)=6$, $K(6)=12$, $K(7)=29$, $K(8)=58$. The bound is attained by survivable sets for every $m$: the extremal configurations of Kleitman's theorem are the balls $\Ball(c,r)$ with $r=\tfr{m-3}2$ for odd $m$ and the doubled balls $\Ball_{m-1}(c,r)\times\{\pm1\}$ of a facet with $r=\tfr{m-4}2$ for even $m$, and these are survivable at the floor ($r=r^\ast(m)$) for odd $m$ and from slack exactly $1$ for even $m$.
\end{corollary}
\begin{proof}
Diameter: a pair at distance $>\delta_m$ is not survivable and survivability is hereditary. For odd $m$, $r^\ast(m)=\lfloor(2m-5)/4\rfloor=\tfr{m-3}2$, so the extremal balls are the floor balls of Corollary~\ref{cor:balls}. For even $m$ take $y$ equal to the centre $c$ on the $m-1$ facet coordinates and $0$ on the doubled coordinate: every cell of the doubled ball has $\langle y,\theta\rangle\ge m-1-2r=3\ge\tfr52$. The matched action of one layer gives the cells of the other layer at distance $r+1$ the drift $2m-5-4(r+1)=-1$, so the first-step bound is $1$; alternating the two layer centres gains at least $3-1=2$ per two steps for every cell with dip $1$, so $\Lstar=1$.
\end{proof}

\begin{proposition}[the instruments are incomparable]\label{prop:incomparable}
Let $T$ be a tetrahedron of the five-cube: four cells pairwise at distance $2$, e.g. $(1,1,1,1,1)$, $(1,1,1,-1,-1)$, $(1,1,-1,1,-1)$, $(1,1,-1,-1,1)$. Every pair of $T$ is survivable (Corollary~\ref{cor:pairs}), every three-subset lies in a radius-one ball, yet $V(T)=2<\tfr52$ with the uniform certificate $q$, because $\sum_\theta\tfr14\theta=(1,1,0,0,0)$. $T$ is never blind-survivable, and from slack $L$ it loses a cell within $L$ steps; by forward search the exact horizon is $2\lfloor(L-1)/2\rfloor$ for $L\le20$, the same as for the far pairs of the five-cube: the two-step alternation that costs one cell $3$ and gains the other $1$ is the slowest way down, and the last two tenths cannot be spent.
\end{proposition}

The proposition answers a question raised about the earlier editions: the ``tetrahedral exclusion'' and the ``averaging bound'' are one statement, and the quantity that decides survivability is not the strength $\sum_i|u_i|=m$ of an instrument---every sign action has the same strength---but the $\ell_1$ norm of the \emph{pooled cell} $\sum_\theta q_\theta\theta$, which measures how much the cells of $S$ agree. The corollaries also give the structure of \emph{minimal obstructions}: by the support bound a non-survivable set contains a non-survivable subset of at most $m+1$ cells, so the minimal non-survivable sets have at most $m+1$ cells; Section~\ref{sec:classification} lists them.

\begin{remark}[one-step duality and sequences]
For a single step, the measure certificate of \cite{Abaee2026minimax} says that no common action keeps $S$ for one period iff some $q\in\Delta(S)$ has $\max_u\sum_\theta q_\theta D(u,\theta)<0$; here that is $2\|\sum q_\theta\theta\|_1-5<0$, the same quantity. The content of Theorem~\ref{thm:criterion} is that over \emph{sequences} with unbounded slack nothing better is possible: time averaging converts any word into a mixed action, and periodic words realise any rational mixed action exactly.
\end{remark}
